\scanpage{68}\mspage{8}
% archive sheet number: 70

par bouchage des trous correspondants reste la $=$ (à isom. can. près). Il faudrait, pour le voir clairement, voir la relation entre la sphère \inserted{\unsure{hol}} $S(I)$ triponctuée, associée à la sphère à bord $\Sigma^*$, et ses paramètres $l_i$ $(i\in I)$ comme déduite \unsure{(précisément)|(rétrécissement)|(réécriture)|(riemanniens)} de $S(I)$ en enlevant des disques \deleted{qui} \inserted{de centre $s_i$?)} $D(l_i)$ \deleted{qui} dont chacun ne dépend que du $l_i$ \unsure{correspondant|correspondants}, pas des autres $l_j$.
\margin{?}

Ce point étant supposé acquis, ou quelque variante, il semble donc bien que (du moins dans le cas hyperbolique orienté) pour une surface conforme \unsure{ponctué|ponctuée} donnée \inserted{de type $(g,\nu)$}, les familles \inserted{de disques $(D_i)$} des disques autour des $a_i$ forme un espace homéom. à $\mathbb R^\nu$ (où $\nu=\mathrm{card}\,I$), i.e. à celui des systèmes d'ombres possibles! Il est alors très tentant de faire la

\uline{Conjecture} Pour \unsure{tte|tt} surface conforme ponctuée, (ou \deleted{\unsure{tout}} au moins dans le cas hyperbolique), l'application \deleted{qui}
\[
(D_i)\longmapsto (o(D_i))
\]
